\documentclass[letterpaper]{article}

\usepackage[margin=1in]{geometry}
\usepackage{amsmath,amssymb}
\usepackage{booktabs}
\usepackage{multirow}
\usepackage{graphicx}
\usepackage{xcolor}
\usepackage{hyperref}
\usepackage[numbers]{natbib}

\title{Token Pruning Catastrophically Fails on Pixel-Shuffle VLAs:\\ Diagnosis and Recovery}

\author{
Anonymous Authors
}

\date{}

\begin{document}
\maketitle

%==============================================================================
\begin{abstract}
%==============================================================================
Inference-time token pruning is a popular strategy for accelerating Vision-Language-Action (VLA) models, yet its interaction with the vision connector architecture is poorly understood.
We present the first systematic study of token pruning across five VLA models (450M--7B parameters) spanning three connector families.
Our central finding is that \textbf{pixel-shuffle connectors are catastrophically incompatible with any form of token pruning}: on SmolVLA (pixel-shuffle, 128 tokens), removing as few as 5\% of vision tokens collapses the success rate from 44.8\% to 0\% on LIBERO manipulation tasks.
In stark contrast, pure linear projection ($\pi$0, 256 tokens) is effectively immune---performance \emph{improves} by 2.4\% even at 90\% token removal.
We propose a zero-cost diagnostic based on the coefficient of variation (CV) of token importance scores: a single forward pass predicts whether a given VLA can tolerate pruning (CV spectrum: 0.038 immune $\rightarrow$ 0.314 catastrophic).
For vulnerable pixel-shuffle architectures, we introduce Vision Token Re-encoding (VTR), a lightweight learned module ($\sim$22M parameters) that recovers success rate from 0\% to 60.1$\pm$1.3\% at 75\% compression, with a flat M-curve confirming near-lossless performance at any compression ratio (M=16: 56.9\% $\approx$ M=128: 58.0\%).
VTR is specifically effective for pixel-shuffle connectors; on MLP-projector VLAs it yields 0\% SR across all compression ratios, confirming the architecture-specificity of both the failure mode and the fix.
We additionally expose a pruning-method confound (L2-norm vs.\ random selection: $\sim$17 percentage point gap on OpenVLA) and high seed sensitivity at architectural transition points (28pp variance on X-VLA at $k$=0.7), both of which complicate cross-model comparisons in prior work.
These findings serve as a cautionary tale for practitioners and provide actionable diagnostic tools for VLA deployment.
\end{abstract}

%==============================================================================
\section{Introduction}
\label{sec:intro}
%==============================================================================

Vision-Language-Action (VLA) models process hundreds of vision tokens per frame, creating a computational bottleneck for real-time robot control.
A growing body of work proposes inference-time token pruning as an efficient solution~\cite{chen2024image,liang2022not}, with several concurrent methods reporting success on specific VLA architectures.
However, these methods are evaluated on models with particular vision connector designs---typically multi-layer MLPs---and the implicit assumption that pruning behavior generalizes across architectures remains untested.

We challenge this assumption with a systematic study across five VLA models spanning three connector families: pure linear projection ($\pi$0), multi-layer MLP (OpenVLA, OpenVLA-OFT), linear projection with deep post-processing (X-VLA with 12-layer BART encoder), and pixel-shuffle (SmolVLA).
Our findings reveal a stark architectural divide.

\textbf{Pixel-shuffle connectors are catastrophically vulnerable.}
SmolVLA's pixel-shuffle connector spatially folds 1024 vision patches into 128 tokens via channel-to-spatial reorganization.
This creates irreversible cross-token entanglement: removing just 5\% of tokens (7 out of 128) annihilates task performance from 44.8\% to 0\% success rate on LIBERO-Object.
No pruning method (random, importance-based, geometric prior) survives.
This finding alone has immediate practical implications---four concurrent works on VLA token pruning all ignore pixel-shuffle architectures entirely.

\textbf{Pure linear projection is immune.}
At the opposite extreme, $\pi$0's pure linear projection produces tokens that are independently mapped from vision patches.
Even removing 90\% of tokens \emph{improves} the action prediction loss by 2.4\%, with importance CV=0.038 (near-perfectly uniform).

Between these extremes, MLP projectors (OpenVLA) and linear-plus-deep-post-processing (X-VLA) show intermediate behaviors, but we present evidence that cross-model comparisons of these intermediate cases are confounded by pruning method choice ($\sim$17pp gap between L2-norm and random pruning on OpenVLA) and high seed sensitivity at transition points (28pp variance on X-VLA at $k$=0.7).
We report these intermediate results as suggestive rather than definitive.

Our contributions are:
\begin{enumerate}
    \item \textbf{Catastrophic failure discovery}: First demonstration that pixel-shuffle vision connectors are incompatible with any form of token pruning, with the failure threshold at $\leq$5\% removal.
    \item \textbf{Zero-cost diagnostic}: A CV-based importance metric that predicts pruning tolerance from a single forward pass (CV: 0.038 immune $\rightarrow$ 0.314 catastrophic).
    \item \textbf{VTR recovery}: A lightweight re-encoding module that recovers pixel-shuffle VLAs from 0\% to 60\%+ SR with a flat compression curve.
    \item \textbf{Confound exposure}: Quantification of pruning-method and seed-sensitivity confounds that complicate cross-architecture comparisons.
\end{enumerate}

%==============================================================================
\section{Background and Related Work}
\label{sec:background}
%==============================================================================

\subsection{Vision-Language-Action Models}

VLAs integrate a pre-trained vision encoder, a vision-to-LLM connector, and a language model backbone for robot control.
The connector transforms vision encoder outputs into token sequences consumed by the LLM.
Connector designs vary substantially:
\textbf{(i)} pure linear projection ($\pi$0~\cite{black2024pi0}: a single linear layer mapping each patch independently);
\textbf{(ii)} multi-layer MLP (OpenVLA~\cite{kim2024openvla}: a 3-layer MLP with GELU activations introducing non-linear cross-channel mixing);
\textbf{(iii)} linear projection with deep post-processing (X-VLA: linear projection followed by a 12-layer BART encoder that introduces cross-token mixing via self-attention);
\textbf{(iv)} pixel-shuffle (SmolVLA~\cite{allal2025smolvla}: spatial folding that reorganizes channel dimensions into spatial dimensions, compressing 1024 patches into 128 tokens).

\subsection{Token Pruning for VLMs and VLAs}

Token pruning reduces inference cost by removing less important vision tokens.
Methods include attention-based importance scoring~\cite{chen2024image}, random subset selection, and geometric priors.
Recent VLA-specific work (Compressor-VLA, SpecPrune-VLA, LightVLA, See What Matters) demonstrates effectiveness on MLP-projector architectures.
Critically, none of these concurrent works evaluates pixel-shuffle VLAs.

%==============================================================================
\section{Method}
\label{sec:method}
%==============================================================================

\subsection{Connector Compressibility Diagnostic}
\label{sec:diagnostic}

We propose a zero-cost diagnostic to predict whether a VLA's vision tokens can be safely pruned.
Given a VLA with $N$ vision tokens, we compute token importance scores $\{s_i\}_{i=1}^{N}$ via a single forward pass using L2-norm of the projected token representations.
The \textbf{coefficient of variation} (CV) of these scores quantifies token importance heterogeneity:
\begin{equation}
    \text{CV} = \frac{\sigma(\{s_i\})}{\mu(\{s_i\})}
\end{equation}
Low CV indicates uniform importance (pruning-tolerant); high CV indicates heterogeneous importance with critical tokens (pruning-vulnerable).
We additionally compute per-layer attention entropy of the LLM's suffix-to-vision attention as a mechanistic supplement:
\begin{equation}
    H_\ell = -\frac{1}{N}\sum_{i=1}^{N} \sum_{j=1}^{N} a_{ij}^{(\ell)} \log a_{ij}^{(\ell)}
\end{equation}
where $a_{ij}^{(\ell)}$ are attention weights at layer $\ell$.

Across our five models, CV forms a continuous spectrum that predicts pruning tolerance:
$\pi$0 (0.038, immune) $<$ X-VLA (0.151, intermediate) $<$ OpenVLA (0.215, graded) $<$ SmolVLA (0.314, catastrophic).

As a practitioner-friendly alternative, we propose a \textbf{lightweight pruning probe}: remove 5\% of tokens and measure action loss change in a single forward pass.
A loss increase $>$20\% signals pruning vulnerability.

\subsection{Vision Token Re-encoding (VTR)}
\label{sec:vtr}

For architectures where pruning fails catastrophically, we propose Vision Token Re-encoding (VTR), a lightweight cross-attention module that \emph{learns} to compress vision tokens rather than deleting them.

Given $N$ input vision tokens $\mathbf{V} \in \mathbb{R}^{N \times d}$, VTR uses $M$ learnable query tokens $\mathbf{Q} \in \mathbb{R}^{M \times d}$ ($M < N$) to produce compressed representations via cross-attention:
\begin{equation}
    \mathbf{Z} = \text{CrossAttn}(\mathbf{Q}, \mathbf{V}, \mathbf{V}) \in \mathbb{R}^{M \times d}
\end{equation}
The module adds $\sim$22M parameters and is trained with the LLM backbone frozen (VTC-only training).

The key insight is that re-encoding \emph{itself}---not compression---is the operative mechanism for pixel-shuffle tokens.
When $M=N$ (passthrough, 0\% compression), VTR still achieves substantial recovery because cross-attention re-organizes the spatially entangled representations into individually meaningful tokens.

%==============================================================================
\section{Experimental Setup}
\label{sec:setup}
%==============================================================================

\textbf{Models.}
We evaluate five VLA models: SmolVLA (450M, pixel-shuffle connector, 128 tokens), X-VLA (880M, linear + 12-layer BART encoder, 50 tokens), $\pi$0 (3.3B, pure linear projection, 256 tokens), OpenVLA (7B, 3-layer MLP with GELU, 256 tokens), and OpenVLA-OFT (7B, same MLP connector, different action head).

\textbf{Benchmark.}
All closed-loop evaluations use LIBERO-Object (10 manipulation tasks, 50 episodes per task, 500 total).
We report overall success rate (SR) with Wilson 95\% confidence intervals.

\textbf{Pruning protocol.}
We evaluate at keep-ratios $k \in \{0.95, 0.9, 0.7, 0.5, 0.3\}$.
For SmolVLA, we test random, importance-based (L2-norm), and geometric (Sobel edge) pruning---all produce identical 0\% results.
For OpenVLA, we evaluate both random pruning (original protocol) and L2-norm pruning to quantify the method confound.
For X-VLA, we use L2-norm pruning.
$\pi$0 is evaluated offline via normalized action-prediction loss.

\textbf{VTR training.}
VTC-only training (frozen LLM backbone) for 100K steps on LIBERO-Object training data.
We sweep $M \in \{16, 32, 64, 96, 128\}$ on SmolVLA and $M \in \{32, 64, 96, 128\}$ on OpenVLA.
All SmolVLA results use 3-seed evaluation (seeds 42, 43, 44).

%==============================================================================
\section{Results}
\label{sec:results}
%==============================================================================

\subsection{Pixel-Shuffle Catastrophic Collapse}
\label{sec:pixel-shuffle}

SmolVLA's pixel-shuffle connector creates extreme vulnerability to token pruning.
Table~\ref{tab:smolvla_cliff} shows that \emph{every} pruning configuration produces 0\% success rate, including removing only 5\% of tokens ($k$=0.95).

\begin{table}[h]
\centering
\caption{SmolVLA pruning results on LIBERO-Object. Baseline SR = 44.8\%. All pruning conditions collapse to 0\%.}
\label{tab:smolvla_cliff}
\begin{tabular}{lccc}
\toprule
\textbf{Keep ratio $k$} & \textbf{Random} & \textbf{L2-norm} & \textbf{Geometric (Sobel)} \\
\midrule
0.95 (5\% removal) & 0\% & 0\% & --- \\
0.9 (10\% removal) & 0\% & 0\% & 0\% \\
0.7 (30\% removal) & 0\% & 0\% & 0\% \\
0.5 (50\% removal) & 0\% & 0\% & 0\% \\
\bottomrule
\end{tabular}
\end{table}

The pixel-shuffle operation folds 1024 SigLIP patches into 128 tokens (16$\times$ compression) by reorganizing spatial dimensions into channels.
Each resulting token encodes a $32 \times 32$ pixel region.
Removing any token eliminates an entire spatial region's information with no redundancy across tokens to compensate.

\subsection{Pure Linear Projection Immunity}
\label{sec:linear}

$\pi$0's pure linear projection produces effectively independent tokens.
At $k$=0.1 (90\% removal), the normalized action-prediction loss is 0.976---a 2.4\% \emph{improvement} over the unpruned model.
The importance CV of 0.038 confirms near-perfect uniformity: all tokens carry approximately equal information, so any subset suffices.
Attention entropy of 0.959 (range [0.931, 0.983] across 18 layers) further confirms that the LLM attends uniformly to vision tokens.

\subsection{Pruning Method Confound: L2-Norm vs.\ Random}
\label{sec:confound}

A critical finding for cross-model comparison: the choice of pruning method introduces a systematic $\sim$17pp confound on MLP-projector VLAs.
Table~\ref{tab:method_confound} shows that L2-norm pruning consistently underperforms random pruning on OpenVLA across all keep-ratios.

\begin{table}[h]
\centering
\caption{OpenVLA pruning method comparison on LIBERO-Object (500 episodes per condition). L2-norm systematically underperforms random by $\sim$17pp.}
\label{tab:method_confound}
\begin{tabular}{lcccc}
\toprule
\textbf{Keep ratio $k$} & \textbf{Random SR (\%)} & \textbf{L2-norm SR (\%)} & \textbf{$\Delta$ (pp)} & \textbf{Significant?} \\
\midrule
0.95 & 69.4 & 52.0 & $-$17.4 & Yes ($>$15pp) \\
0.90 & 62.0 & 44.0 & $-$18.0 & Yes ($>$15pp) \\
0.70 & 49.0 & 31.0 & $-$18.0 & Yes ($>$15pp) \\
0.50 & 37.4 & 21.0 & $-$16.4 & Yes ($>$15pp) \\
0.30 & 20.6 & 7.0 & $-$13.6 & Marginal \\
\bottomrule
\end{tabular}
\end{table}

The L2-norm importance metric is a misleading proxy for MLP-projected tokens.
The 3-layer MLP with GELU activations compresses action-relevant information into low-magnitude tokens: after non-linear transformation, tokens with small L2-norm can carry high gradient-based importance (Spearman $\rho$ = 0.407 between L2-norm and gradient importance).
L2-norm pruning systematically removes these task-critical, low-magnitude tokens.
Random pruning, by contrast, samples uniformly across importance quartiles.

This confound has direct implications for the literature: if Model A is evaluated with random pruning and Model B with importance-based pruning, the observed gap conflates architectural differences with method artifacts.
Both 95\% confidence intervals (Wilson) are non-overlapping at every $k$-value, confirming this is a systematic effect, not noise.

\textbf{Cross-seed reproducibility of random pruning.}
To confirm that the random pruning baseline itself is stable, we evaluate OpenVLA random pruning with an alternative seed (seed43 vs.\ original seed7).
At $k$=0.7: 46.0\% [41.7, 50.4] vs.\ original 49.0\% ($\Delta$=$-$3.0pp).
At $k$=0.5: 35.0\% [31.0, 39.2] vs.\ original 37.4\% ($\Delta$=$-$2.4pp).
Both deltas fall within the Wilson CI, confirming that the $\sim$17pp L2-norm vs.\ random gap is a method effect, not seed variance.

\subsection{Seed Sensitivity at Architectural Transition Points}
\label{sec:seed}

X-VLA exhibits a sharp performance cliff between $k$=0.9 (97.6\% SR) and $k$=0.7 (56.2\% SR)---a 41pp drop.
We test the stability of this transition by evaluating with an alternative seed (seed44 vs.\ original seed42).

\begin{table}[h]
\centering
\caption{X-VLA cross-seed comparison on LIBERO-Object (500 episodes).}
\label{tab:xvla_seed}
\begin{tabular}{lccc}
\toprule
\textbf{Metric} & \textbf{Seed 42} & \textbf{Seed 44} & \textbf{$\Delta$} \\
\midrule
$k$=0.9 SR & 97.6\% & 91.0\% [88.3, 93.2] & $-$6.6pp \\
$k$=0.7 SR & 56.2\% & 28.0\% [24.2, 32.1] & $-$28.2pp \\
Tasks $>$60\% & 6 & 2 & $-$4 \\
Tasks $<$10\% & 3 & 4 & $+$1 \\
\bottomrule
\end{tabular}
\end{table}

At $k$=0.7 (the transition point), seed44 produces a 28.2pp lower SR than seed42---far exceeding typical 500-episode variance ($<$8pp).
This indicates that the cliff position itself is seed-sensitive, spanning approximately $k \in [0.7, 0.9]$.
At $k$=0.9 (above the cliff), seed variance is modest (6.6pp), and both seeds remain above 90\%.

This has two implications: (1) single-seed evaluations at architectural transition points can be highly misleading, and (2) characterizing X-VLA's cliff as occurring at precisely $k$=0.7 overstates the precision of the measurement.
We can only reliably claim that X-VLA exhibits a cliff ``somewhere in $k \in [0.7, 0.9]$.''

\subsection{VTR Recovers Pixel-Shuffle VLAs}
\label{sec:vtr_results}

VTR successfully recovers SmolVLA from the catastrophic 0\% pruning cliff.
Table~\ref{tab:vtr_mcurve} shows the full M-sweep with 3-seed evaluation.

\begin{table}[h]
\centering
\caption{SmolVLA VTR M-curve on LIBERO-Object (3-seed mean $\pm$ std). Baseline SR = 44.8\%, hard pruning SR = 0\% at all $k$. All M-values recover to $>$54\%.}
\label{tab:vtr_mcurve}
\begin{tabular}{lcccc}
\toprule
\textbf{$M$} & \textbf{Compression} & \textbf{Mean SR (\%)} & \textbf{Std} & \textbf{vs.\ Pruning} \\
\midrule
16 & 87.5\% & 56.9 & $\pm$1.3 & $+$56.9pp \\
32 & 75\% & 60.1 & $\pm$1.3 & $+$60.1pp \\
64 & 50\% & 62.0 & $\pm$2.3 & $+$62.0pp \\
96 & 25\% & 54.5 & $\pm$3.2 & $+$54.5pp \\
128 & 0\% (passthrough) & 58.0 & $\pm$3.5 & $+$58.0pp \\
\bottomrule
\end{tabular}
\end{table}

The M-curve is flat: all values fall within a 7.5pp range (54.5--62.0\%), with M=16 (87.5\% compression) statistically indistinguishable from M=128 (passthrough, $\Delta$=1.1pp).
This confirms that VTR achieves near-lossless compression at any ratio for pixel-shuffle connectors.
The $M$=128 passthrough result is particularly informative: with zero compression, VTR still recovers from 0\% to 58.0\%, proving that the operative mechanism is \emph{re-encoding} (disentangling spatially entangled tokens) rather than compression.

\textbf{Architecture independence of re-encoding.}
An MLP adapter with matched parameter count ($\sim$22.14M) achieves 63.4$\pm$2.6\% at $M$=32 vs.\ cross-attention's 60.1$\pm$1.3\% ($p$=0.15, not significant).
Both architectures recover from 0\% and exceed the uncompressed baseline, confirming that re-encoding---not a specific architecture---is the causal mechanism.

\textbf{Cross-suite generalization.}
VTR trained on LIBERO-Object transfers to other suites: $M$=32 achieves 53.2\% on Spatial (vs.\ 0.6\% pruning) and 56.4\% on Goal (vs.\ 0\% pruning), demonstrating robust generalization from the catastrophic cliff.

\subsection{VTR Fails on MLP-Projector VLAs}
\label{sec:vtr_mlp}

VTR's effectiveness is architecture-specific.
Table~\ref{tab:vtr_openvla} shows that VTC-only training on OpenVLA (MLP projector) produces 0\% SR across all M-values.

\begin{table}[h]
\centering
\caption{OpenVLA VTR results (VTC-only, 100K steps, LIBERO-Object). All M-values fail.}
\label{tab:vtr_openvla}
\begin{tabular}{lccl}
\toprule
\textbf{$M$} & \textbf{Compression} & \textbf{SR (\%)} & \textbf{Failure mode} \\
\midrule
32 & 87.5\% & 0 & Action collapse (identical output) \\
64 & 50\% & 2 & Near-collapse (sporadic) \\
96 & 25\% & 0 & Action collapse \\
128 & 0\% (passthrough) & 0 & Action collapse at 20K steps \\
\bottomrule
\end{tabular}
\end{table}

The failure mode is action collapse: the model outputs identical actions every timestep regardless of observation.
Training loss converges (0.028 at 100K steps) but eval loss diverges (0.14), indicating severe overfitting to the training distribution.

This result has a clear mechanistic explanation.
MLP-projected tokens are \emph{already individually meaningful}---each token independently encodes its corresponding vision patch through the 3-layer MLP.
VTR's cross-attention restructures this representation, breaking the per-token independence that the frozen LLM backbone relies upon.
For pixel-shuffle tokens (spatially entangled, CV=0.314), VTR's re-encoding is beneficial because it \emph{disentangles} the spatial folding.
For MLP-projected tokens (partially independent, CV=0.215), re-encoding is harmful because it \emph{destroys} existing structure.

$M$=64 achieves a marginal 2\% SR (10/500), but the 95\% CI [1.0, 3.8] overlaps with chance-level performance, and we do not consider this meaningful recovery.
The universal failure across $M \in \{32, 64, 96, 128\}$ rules out compression ratio as the explanation---even passthrough ($M$=128, 0\% compression) fails, confirming that the VTR architecture itself is incompatible with MLP-projected tokens.

\subsection{Diagnostic Metric Validation}
\label{sec:diagnostic_results}

Table~\ref{tab:diagnostic} summarizes the CV-based diagnostic across all five models, confirming the continuous predictive spectrum.

\begin{table}[h]
\centering
\caption{Connector compressibility diagnostic. CV predicts pruning tolerance along a continuous spectrum.}
\label{tab:diagnostic}
\begin{tabular}{llccl}
\toprule
\textbf{Model} & \textbf{Connector} & \textbf{CV} & \textbf{Entropy} & \textbf{Pruning behavior} \\
\midrule
$\pi$0 (3.3B) & Pure linear & 0.038 & 0.959 & Immune ($k$=0.1: $+$2.4\%) \\
X-VLA (880M) & Linear + BART 12L & 0.151 & 0.614 & Cliff in $k \in [0.7, 0.9]$ \\
OpenVLA (7B) & 3-layer MLP & 0.215 & 0.661 & Graded ($k$=0.5: $-$49\%) \\
OFT (7B) & 3-layer MLP & --- & 0.667 & $\approx$ OpenVLA (entropy match) \\
SmolVLA (450M) & Pixel-shuffle & 0.314 & 0.952 & Catastrophic ($k$=0.95: 0\%) \\
\bottomrule
\end{tabular}
\end{table}

CV captures the degree of non-linear mixing in the vision-to-LLM pathway: pure linear projection produces near-zero CV (no mixing, tokens independent), while pixel-shuffle produces the highest CV (maximal spatial entanglement).
The two intermediate architectures (X-VLA with BART post-processing and OpenVLA with MLP) produce intermediate CVs that correctly predict intermediate pruning sensitivity.

The diagnostic is practitioner-friendly: it requires a single forward pass, no task-specific evaluation, and produces a scalar that directly indicates whether pruning is safe.
CV $<$ 0.05 suggests pruning immunity; CV $>$ 0.25 signals catastrophic risk.

%==============================================================================
\section{Discussion}
\label{sec:discussion}
%==============================================================================

\subsection{Implications for Practitioners}

Our results provide three actionable guidelines for VLA deployment:

\textbf{(1) Never prune pixel-shuffle VLAs.}
The catastrophic failure at 5\% removal leaves no safety margin.
If a VLA uses pixel-shuffle (or any spatial folding connector), practitioners must use learned compression (VTR) or avoid token reduction entirely.

\textbf{(2) Use the CV diagnostic before deploying pruning.}
A single forward pass determines whether a given VLA can tolerate pruning.
This is especially important as new VLA architectures proliferate---the connector design may not be immediately obvious from the model card.

\textbf{(3) Match the compression strategy to the connector.}
VTR is effective for pixel-shuffle ($0\% \rightarrow 60\%+$) but harmful for MLP projectors (0\% across all M).
Simple pruning works for linear projection.
There is no universal token reduction strategy.

\subsection{Why Four Concurrent Works Missed This}

Compressor-VLA, SpecPrune-VLA, LightVLA, and See What Matters all evaluate on MLP-projector VLAs (primarily OpenVLA-family models) where pruning degrades gracefully.
None tests pixel-shuffle architectures.
Our finding that pixel-shuffle VLAs catastrophically fail under any pruning is invisible to evaluations restricted to MLP-projector models.
As pixel-shuffle connectors are used in production VLAs (SmolVLA), this is not merely a theoretical concern.

\subsection{Limitations}

\textbf{Single benchmark.}
All closed-loop evaluations use LIBERO simulation.
Real-robot validation and additional benchmarks (e.g., CALVIN, RLBench) are needed to confirm generality.

\textbf{Intermediate classes are suggestive.}
The MLP-graded (OpenVLA) and threshold-graded (X-VLA) categories are confounded by pruning method differences ($\sim$17pp L2-norm vs.\ random gap) and seed sensitivity (28pp at $k$=0.7).
We present these as suggestive evidence of intermediate behavior, not as definitive architectural characterizations.
Clean cross-model comparison requires standardized pruning protocols.

\textbf{Offline evaluation for $\pi$0.}
$\pi$0 pruning tolerance is measured via offline action-prediction loss on pre-collected data, not closed-loop evaluation.
While the results are striking (loss improvement at 90\% removal), closed-loop confirmation would strengthen the claim.

\textbf{Perceiver-style connectors untested.}
We do not evaluate Perceiver Resampler connectors (e.g., RoboFlamingo).
These use learned query tokens to cross-attend to vision features and represent a fourth connector family whose pruning behavior is unknown.
We leave this to future work.

\textbf{VTR training regime.}
VTR is trained with VTC-only (frozen backbone).
Joint fine-tuning may yield different results, particularly for MLP-projector models where VTC-only training fails.

%==============================================================================
\section{Conclusion}
\label{sec:conclusion}
%==============================================================================

We present the first systematic study of token pruning across VLA vision connector architectures.
Our central finding is that pixel-shuffle connectors are catastrophically incompatible with token pruning: removing as few as 5\% of vision tokens collapses task performance to 0\%.
This stands in stark contrast to pure linear projection, which is immune to even 90\% token removal.
We provide a zero-cost CV-based diagnostic that predicts pruning tolerance from a single forward pass, and Vision Token Re-encoding (VTR), a lightweight module that recovers pixel-shuffle VLAs from 0\% to 60\%+ SR at any compression ratio.
Critically, VTR is a pixel-shuffle-specific solution: it fails on MLP-projector VLAs, reinforcing the architecture-dependency of both the problem and the fix.
We additionally expose pruning-method and seed-sensitivity confounds that caution against naive cross-architecture comparisons.

These results serve as both a cautionary tale---practitioners must verify connector compatibility before deploying token pruning---and a practical toolkit: the CV diagnostic and VTR together provide a diagnose-then-compress framework for safe VLA token reduction.

%==============================================================================
% References
%==============================================================================
\bibliographystyle{plainnat}

\begin{thebibliography}{10}

\bibitem{kim2024openvla}
M.~J. Kim, K.~Pertsch, S.~Karamcheti, T.~Xiao, A.~Balakrishna, S.~Nair, R.~Rafailov, E.~Foster, G.~Lam, P.~Sanketi, Q.~Vuong, T.~Kollar, B.~Burchfiel, R.~Tedrake, D.~Sadigh, S.~Levine, P.~Liang, and C.~Finn.
\newblock Open{VLA}: An open-source vision-language-action model.
\newblock In \emph{CoRL}, 2024.

\bibitem{allal2025smolvla}
L.~B. Allal, F.~Mazzotta, A.~Luccioni, and T.~Wolf.
\newblock Smol{VLA}: A vision-language-action model for affordable robots.
\newblock \emph{arXiv preprint}, 2025.

\bibitem{black2024pi0}
K.~Black, N.~Brown, D.~Driess, A.~Esmail, M.~Equi, C.~Finn, N.~Fusai, L.~Groom, K.~Hausman, B.~Ichter, S.~Jakber, T.~Jones, L.~Ke, S.~Levine, A.~Li, M.~Liwag, A.~Majumdar, J.~Malik, A.~Rao, D.~Sadigh, P.~Sanketi, J.~Song, A.~Toshev, S.~Vuong, Z.~Wang, S.~Welker, and C.~Xu.
\newblock $\pi_0$: A vision-language-action flow model for general robot control.
\newblock \emph{arXiv preprint arXiv:2410.24164}, 2024.

\bibitem{chen2024image}
L.~Chen, J.~Zhao, C.~Zhuo, J.~Chen, and Z.~Lu.
\newblock An image is worth 1/2 tokens after layer 2: Plug-and-play inference acceleration for large vision-language models.
\newblock In \emph{ECCV}, 2024.

\bibitem{liang2022not}
Y.~Liang, C.~Ge, Z.~Tong, Y.~Song, J.~Wang, and P.~Xie.
\newblock Not all patches are what you need: Expediting vision transformers via token reorganizations.
\newblock In \emph{ICLR}, 2022.

\end{thebibliography}

\end{document}
